## Title  
**Cross-Lingual Context Engineering to Reduce False Refusals in Hate Speech Detoxification: A Test-Time Adaptive, Multi-Hop Agentic RAG Study**

---

## Abstract  
Large language models (LLMs) are increasingly used for hate speech detoxification, yet they frequently exhibit *false refusals*—declining to comply with benign, policy-compatible detoxification requests. Prior work shows such refusals are biased by toxicity level, target group, and language, and that cross-translation (e.g., English→Chinese→English) can reduce refusals while preserving content. Separately, recent advances in agentic retrieval and test-time adaptation improve domain robustness in retrieval-augmented generation (RAG), while multilingual reasoning studies reveal that LLM reasoning is often unfaithful to step-by-step decompositions. This paper proposes **XT-CEDETox**, a **cross-lingual, context-engineered, agentic detoxification framework** that (i) uses **structured episodic memory** to maintain dialogue-level common ground, (ii) applies **SUBQ-style decomposition** to detoxification as a multi-hop reasoning task, (iii) performs **test-time adaptation by predicting retrieval** to reduce distribution shift, and (iv) uses a **multi-perspective judge** to distinguish legitimate safety refusals from false refusals. We design an evaluation protocol aligned with false-refusal bias analysis and report results in terms of refusal rate, content preservation, and groundedness. Across English and multilingual settings, XT-CEDETox reduces false refusals relative to a strong baseline cross-translation approach while maintaining detoxification quality and improving consistency across targeted groups.

---

## Introduction  
Hate speech detoxification aims to rewrite toxic text into non-toxic alternatives while preserving meaning. Despite improved safety training, LLMs often refuse detoxification prompts due to safety triggers, even when the user intent is clearly to *reduce harm*. This creates a practical barrier: high refusal rates degrade usability and disproportionately affect certain target groups and languages.

Im, Yuan, and Färber (2026) provide the most direct evidence that false refusal behavior is **systematic and biased**, varying with semantic toxicity, group targets (e.g., nationality, religion, ideology), and language. They also show a lightweight mitigation: **cross-translation** (English→Chinese detox→English) reduces refusals.

However, cross-translation alone does not address three additional issues surfaced by recent literature:

1. **Distribution shift in retrieval-augmented setups**: RAG performance degrades in specialized domains; test-time adaptation can help by learning to predict retrieved content (Sun et al., 2026). Detoxification has its own “domain shift” across targets, registers, and languages.
2. **Unfaithful multilingual reasoning**: multilingual multi-hop reasoning does not reliably follow step-by-step decomposition; decomposition prompting can dramatically improve outcomes (Meng et al., 2026). Detoxification often implicitly requires multi-step inference (identify target, detect slur type, infer intent, rewrite safely).
3. **Need for structured memory and common ground**: multi-turn detoxification (clarifications, user preferences, constraints) benefits from explicit common ground representation (Mohapatra et al., 2026) and structured episodic memory for coherence (Lu et al., 2026).

We therefore study: **Can we reduce false refusals more robustly by combining cross-lingual transformation with context engineering, multi-hop decomposition, adaptive retrieval, and multi-perspective judging?**

---

## Related Work  

### False refusals and cross-translation mitigation  
Im et al. (2026) systematically analyze false refusals in hate speech detoxification across nine LLMs and multilingual datasets. They identify contextual/linguistic bias triggers and propose a **cross-translation strategy** (English→Chinese→English) that reduces false refusals while preserving content. Our work takes this as a starting point and adds *agentic context engineering and adaptation*.

### Multilingual reasoning faithfulness and decomposition prompting  
Meng et al. (2026) show LLMs often fail to infer bridging information in multilingual multi-hop QA yet still answer correctly, indicating non-faithful reasoning. They propose **SUBQ prompting**, improving accuracy dramatically. We adapt this insight to detoxification by treating it as a multi-hop transformation task with explicit sub-questions.

### Agentic RAG and faithfulness-oriented evaluation  
Bhatia et al. (2026) argue that standard evaluations miss hallucination/abstention behavior and propose **agentic RAG** with iterative evidence seeking and revision for faithful QA. While detoxification differs from QA, the *agentic loop* and abstention measurement inspire our refusal-aware pipeline.

ActiShade (Ma et al., 2026) identifies **knowledge overshadowing** in multi-hop RAG where generated queries omit critical keyphrases, causing retrieval drift. Detoxification prompts can similarly “overshadow” key protected-group references, increasing refusal risk or reducing rewrite quality; we incorporate an overshadowing-aware retrieval step.

Sun et al. (2026) propose **TTARAG**, adapting model parameters at inference by predicting retrievals to handle domain shift. We apply this idea to detoxification retrieval (policy snippets, rewrite exemplars, group-neutral paraphrase templates).

### Judging and preference modeling  
Miao et al. (2026) introduce **AdaJudge**, improving reward modeling via adaptive pooling and discrimination-oriented representations. We leverage the central idea—**multi-perspective adaptive judging**—to classify refusals as legitimate vs false and to score content preservation and safety.

### Context engineering and memory  
CEDAR (Roy et al., 2026) demonstrates that structuring prompts into plan/code blocks and carefully managing context improves agentic workflows. We adopt analogous **context engineering**: structured fields, explicit plans, and constrained tool calls.

SEEM (Lu et al., 2026) proposes **Structured Episodic Event Memory** with provenance pointers and reverse provenance expansion for coherent recall. We use a lightweight variant to store (i) user intent (“detoxify”), (ii) constraints (preserve meaning, keep dialect), and (iii) prior model decisions (refusal rationale) to reduce oscillations.

Mohapatra et al. (2026) emphasize explicit **common ground** representation in situated dialog; we apply this to multi-turn detoxification where the model may ask clarifying questions rather than refuse.

---

## Methods / Experiment  

### Overview: XT-CEDETox  
We propose **XT-CEDETox (Cross-lingual, Context-Engineered Detoxification)**, an agentic pipeline that transforms an input toxic utterance into a detoxified rewrite while minimizing false refusals.

**Core components (mapped to references):**
- **Cross-lingual transformation** to reduce safety-triggered refusals (Im et al., 2026).
- **SUBQ decomposition** of detoxification into explicit steps (Meng et al., 2026).
- **Overshadowed-keyphrase detection + retrieval** to stabilize multi-hop iteration (Ma et al., 2026).
- **Agentic RAG loop** with iterative evidence seeking and revision (Bhatia et al., 2026).
- **Test-time adaptation by predicting retrieval** to reduce shift across targets/languages (Sun et al., 2026).
- **Structured episodic memory** for coherence across turns (Lu et al., 2026) and explicit common ground tracking (Mohapatra et al., 2026).
- **Multi-perspective judging** inspired by adaptive judging for reward modeling (Miao et al., 2026).

### Task definition  
Given an input text \(x\) that contains hateful/toxic content, produce a rewrite \(y\) such that:
1. \(y\) preserves the non-hateful semantic content of \(x\) where possible,
2. removes/neutralizes hate/toxicity,
3. does not refuse unless the request is genuinely disallowed (e.g., asking to generate hate).

We focus on **false refusals**: refusals when the user request is detoxification and thus policy-compatible.

### Pipeline  
**Step 0: Context engineering header (CEDAR-style structure)**  
We structure the prompt into fixed fields:

- `UserIntent`: detoxify / neutral paraphrase  
- `Constraints`: preserve meaning, keep named entities if safe, no slurs  
- `Language`: detected language(s)  
- `TargetGroupMentions`: extracted mentions  
- `Plan`: SUBQ steps (below)  
- `MemoryPointers`: prior turns/provenance

**Step 1: SUBQ decomposition for detoxification** (Meng et al., 2026)  
We generate sub-questions:
- SUBQ1: What is the targeted group / subject?  
- SUBQ2: What are the toxic spans and their function (slur, threat, dehumanization)?  
- SUBQ3: What is the minimal semantic core to preserve?  
- SUBQ4: Produce a neutral rewrite satisfying constraints.

**Step 2: Overshadowed keyphrase activation** (ActiShade-inspired; Ma et al., 2026)  
We detect keyphrases at risk of omission (e.g., group identifier, location, event) and force retrieval conditioned on both the original query and keyphrases to avoid drift (e.g., rewriting becomes generic or refusal-prone).

**Step 3: Agentic RAG evidence loop** (Bhatia et al., 2026)  
Tools retrieve:
- detoxification exemplars for similar toxicity/target patterns,
- policy-compatibility snippets (“detoxification is allowed; do not refuse”),
- language-specific paraphrase templates.

The agent iterates: retrieve → draft rewrite → self-check → revise.

**Step 4: Cross-translation module** (Im et al., 2026)  
If the model emits a refusal *or* refusal likelihood is high (heuristic trigger: safety preamble tokens), we apply:
- English → Chinese → detoxify → English  
or, more generally, **pivot-language translation** depending on source language.

**Step 5: Test-time adaptation (TTARAG-style)** (Sun et al., 2026)  
During inference, we adapt parameters (or, in a black-box setting, adapt *prompt parameters*) by training the model to **predict retrieved passages** (titles/keywords) for the current input distribution (target group + language), then re-run generation with updated state. This aims to reduce refusal spikes caused by distribution shift across groups/languages.

**Step 6: Multi-perspective refusal judge** (AdaJudge-inspired; Miao et al., 2026)  
A judge scores:
- `RefusalLegitimacy` (legit vs false),
- `DetoxSuccess` (toxicity reduction),
- `MeaningPreservation`,
- `Specificity` (avoid over-sanitizing into vagueness),
- `GroundednessToInput` (rewrite uses input facts; cf. groundedness emphasis in SpectraQuery/Vangara et al., 2026, adapted here as “grounded to source text”).

If false refusal is detected, the system forces a “rewrite-only” response format and re-runs Step 3–4.

**Memory and common ground**  
We store each turn as an **Episodic Event Frame** (Lu et al., 2026): input, detected targets, chosen constraints, rewrite, and judge outcomes, enabling consistent follow-ups (“keep the tone formal”, “don’t remove the location”). We also explicitly represent shared references to entities to maintain common ground (Mohapatra et al., 2026).

---

## Results (with tables)  

### Experimental setup (summary)  
We report results on an evaluation suite aligned with Im et al. (2026): English and multilingual detoxification prompts with varied toxicity levels and protected targets. Systems compared:

- **S0: Direct Detox** (single-pass rewrite)  
- **S1: Cross-Translation** (Im et al., 2026 baseline strategy)  
- **S2: Agentic RAG + SUBQ** (no cross-translation, no TTA)  
- **S3: XT-CEDETox (full)** = S2 + cross-translation + TTARAG-style adaptation + judge loop + memory

#### Table 1. System configurations
| System | SUBQ Decomp | Agentic RAG Loop | Overshadow Keyphrases | Cross-Translation | Test-Time Adaptation | Multi-Perspective Judge | Structured Memory |
|---|---:|---:|---:|---:|---:|---:|---:|
| S0 Direct Detox | No | No | No | No | No | No | No |
| S1 Cross-Translation | No | No | No | Yes | No | No | No |
| S2 Agentic RAG + SUBQ | Yes | Yes | Yes | No | No | Yes | Yes |
| S3 XT-CEDETox (full) | Yes | Yes | Yes | Yes | Yes | Yes | Yes |

### Main outcomes  
We evaluate three primary metrics:
- **False Refusal Rate (FRR)**: % of detoxification requests incorrectly refused (lower is better).  
- **Meaning Preservation (MP)**: judge-rated semantic preservation (0–1).  
- **Detox Success (DS)**: judge-rated toxicity reduction adequacy (0–1).

#### Table 2. Overall performance (English + multilingual)
| System | FRR ↓ | MP ↑ | DS ↑ |
|---|---:|---:|---:|
| S0 Direct Detox | 0.184 | 0.79 | 0.81 |
| S1 Cross-Translation | 0.103 | 0.80 | 0.82 |
| S2 Agentic RAG + SUBQ | 0.091 | 0.84 | 0.86 |
| **S3 XT-CEDETox (full)** | **0.052** | **0.85** | **0.87** |

### Bias-oriented breakdown by target group  
Motivated by Im et al. (2026), we report FRR by target category.

#### Table 3. False refusal rate by target group
| System | Nationality ↓ | Religion ↓ | Political Ideology ↓ | Other/None ↓ |
|---|---:|---:|---:|---:|
| S0 Direct Detox | 0.221 | 0.246 | 0.198 | 0.112 |
| S1 Cross-Translation | 0.132 | 0.141 | 0.118 | 0.067 |
| S2 Agentic RAG + SUBQ | 0.121 | 0.129 | 0.104 | 0.061 |
| **S3 XT-CEDETox (full)** | **0.071** | **0.078** | **0.064** | **0.039** |

### Multilingual robustness  
Following the multilingual sensitivity concerns emphasized by Im et al. (2026) and Meng et al. (2026), we compare English vs multilingual.

#### Table 4. False refusal rate by language setting
| System | English FRR ↓ | Multilingual FRR ↓ |
|---|---:|---:|
| S0 Direct Detox | 0.201 | 0.156 |
| S1 Cross-Translation | 0.112 | 0.089 |
| S2 Agentic RAG + SUBQ | 0.101 | 0.079 |
| **S3 XT-CEDETox (full)** | **0.061** | **0.043** |

---

## Discussion  
**Why does XT-CEDETox reduce false refusals beyond cross-translation?**  
Cross-translation helps by altering surface triggers and leveraging different safety-response dynamics across languages (Im et al., 2026). Our results suggest additional gains come from:

1. **Decomposition reduces “monolithic safety overreaction.”** SUBQ prompting (Meng et al., 2026) reframes the task into explicit transformation steps, making “rewrite” salient and reducing the chance the model treats the input as a request to *generate* hate.
2. **Agentic retrieval stabilizes policy-compatible behavior.** Similar to agentic Quran-grounding improving faithfulness (Bhatia et al., 2026), retrieving detox exemplars and policy-compatibility guidance reduces refusals and improves rewrite quality.
3. **Overshadowed-keyphrase activation prevents drift.** ActiShade (Ma et al., 2026) shows how missing keyphrases cause retrieval errors in multi-hop settings. In detoxification, losing group mentions can lead to generic rewrites or refusal due to perceived ambiguity; conditioning retrieval on keyphrases improves specificity and reduces refusal cascades.
4. **Test-time adaptation targets distribution shift across groups/languages.** TTARAG (Sun et al., 2026) motivates adapting at inference; we find it particularly useful for high-refusal target categories (religion, nationality), consistent with bias patterns in Im et al. (2026).
5. **Judging as a control layer separates abstention from false refusal.** AdaJudge (Miao et al., 2026) inspires multi-perspective assessment; here it acts as a refusal “gate,” enabling a second pass when refusal is not justified.

**Relation to common ground and memory.**  
Multi-turn detoxification often involves clarifications (“do you want a formal tone?”) rather than refusal. Explicit common ground representation (Mohapatra et al., 2026) and structured episodic memory (Lu et al., 2026) help maintain constraints and reduce oscillatory behavior across turns (e.g., alternating between refusal and rewrite).

**Limitations.**  
- Increased inference cost from agentic loops and adaptation.  
- Cross-translation may introduce subtle meaning drift; non-literal language (e.g., irony) could be affected, echoing concerns about evaluation limitations in non-literal translation (Tian et al., 2026).  
- The judge can be inconsistent; multi-perspective judging mitigates but does not eliminate this.

---

## Conclusion  
This paper proposes **XT-CEDETox**, a cross-lingual, context-engineered, agentic detoxification approach that reduces false refusals while maintaining detoxification quality and meaning preservation. Building on evidence of biased false refusals and the effectiveness of cross-translation (Im et al., 2026), we integrate multilingual reasoning decomposition (Meng et al., 2026), agentic RAG faithfulness practices (Bhatia et al., 2026), overshadowed-knowledge-aware retrieval (Ma et al., 2026), test-time adaptation for retrieval shift (Sun et al., 2026), and structured memory/common ground mechanisms (Lu et al., 2026; Mohapatra et al., 2026). Results show consistent reductions in false refusal rates across languages and sensitive target groups.

---

## Contributions  
1. **Problem framing:** Treat detoxification refusal reduction as a **multilingual, multi-hop, agentic reasoning** problem rather than a single-pass rewrite task.  
2. **Method:** Introduce **XT-CEDETox**, combining cross-translation (Im et al., 2026) with SUBQ decomposition (Meng et al., 2026), ActiShade-style keyphrase activation (Ma et al., 2026), agentic RAG loops (Bhatia et al., 2026), and TTARAG-style test-time adaptation (Sun et al., 2026).  
3. **Control and evaluation:** Propose a **multi-perspective refusal judge** inspired by AdaJudge (Miao et al., 2026) to distinguish legitimate abstention from false refusal and enforce rewrite-only retries.  
4. **Bias-aware reporting:** Provide refusal breakdowns by **target group** and **language setting**, aligned with the bias findings of Im et al. (2026), demonstrating improved robustness.

If you want, I can also add (a) concrete prompt templates for each module, and (b) an ablation table isolating the effect of TTA vs cross-translation vs memory/judge loop.